% !TEX program = platex
% 数学ノート集・共通テンプレート v2（積分論2026年度第1部の版面を参考）
% A5判。章内で定義・定理・例・注意等の番号を共通化。
% コンパイル： platex -interaction=nonstopmode math_note_template_v2.tex (2回)
%              dvipdfmx math_note_template_v2.dvi
\documentclass[a5paper,10pt,onecolumn]{jsarticle}
\usepackage[dvipdfmx,top=22mm,bottom=20mm,inner=19mm,outer=18mm,headheight=13pt,headsep=6mm]{geometry}
\usepackage{amsmath,amssymb,amsthm,mathtools}
\usepackage{enumitem,fancyhdr,titlesec}
\usepackage[dvipdfmx,bookmarks=true,colorlinks=false]{hyperref}
\usepackage{pxjahyper}

% ----- 差し替える文書情報 -----
\newcommand{\NoteTitle}{97.050 二次形式と極値}
\newcommand{\NoteSeries}{数学ノート集}

% ----- 文字・組版 -----
% 参考PDFの質素な明朝＋Computer Modernの印象を引き継ぎながら，
% A5判で窮屈にならないよう本文の行間を少し確保する．
\linespread{1.13}
\setlength{\parindent}{1zw}
\setlength{\parskip}{0pt}
\setlength{\abovedisplayskip}{6pt plus 2pt minus 2pt}
\setlength{\belowdisplayskip}{6pt plus 2pt minus 2pt}
\setlength{\abovedisplayshortskip}{3pt plus 2pt minus 1pt}
\setlength{\belowdisplayshortskip}{4pt plus 2pt minus 1pt}
\setlist[enumerate]{topsep=3pt,itemsep=1pt,parsep=0pt,leftmargin=2.4zw}
\titleformat{\section}{\normalfont\normalsize\bfseries}{\thesection}{0.8zw}{}
\titlespacing*{\section}{0pt}{2.0ex plus .4ex}{.9ex}
\titleformat{\subsection}{\normalfont\normalsize\bfseries}{\thesubsection}{0.7zw}{}
\titlespacing*{\subsection}{0pt}{1.5ex plus .3ex}{.6ex}
\raggedbottom

% ----- 柱とページ番号：ページ上部、罫線なし -----
\fancypagestyle{mathnote}{%
  \fancyhf{}%
  \fancyhead[L]{\footnotesize \NoteSeries\ （\NoteTitle）}%
  \fancyhead[R]{\footnotesize \thepage}%
  \renewcommand{\headrulewidth}{0pt}%
  \renewcommand{\footrulewidth}{0pt}%
}
\pagestyle{mathnote}

% ----- 節内通し番号の定義・定理群 -----
\newtheoremstyle{mathnotetheorem}
 {4.5pt}{4.5pt}{\normalfont}{}%
 {\bfseries}{}{0.65zw}%
 {\thmname{#1}\thmnumber{ #2}\thmnote{（#3）}}
\theoremstyle{mathnotetheorem}
\newtheorem{theorem}{定理}[section]
\newtheorem{definition}[theorem]{定義}
\newtheorem{lemma}[theorem]{補題}
\newtheorem{proposition}[theorem]{命題}
\newtheorem{corollary}[theorem]{系}
\newtheorem{example}[theorem]{例}
\newtheorem{remark}[theorem]{注意}
\renewcommand{\proofname}{証明}
\renewcommand{\qedsymbol}{$\square$}
\makeatletter
\renewenvironment{proof}[1][\proofname]{%
  \par\pushQED{\qed}\normalfont\topsep4\p@\@plus2\p@\relax
  \trivlist\item[\hskip\labelsep\bfseries #1]\ignorespaces
}{\popQED\endtrivlist\@endpefalse}
\makeatother

% ----- 使うノートで適宜拡張する記号 -----
\newcommand{\R}{\mathbb{R}}
\newcommand{\norm}[1]{\lVert #1\rVert}
\newcommand{\abs}[1]{\lvert #1\rvert}

\begin{document}
\section{入試の最大値問題を二次形式へ変換する}
本ノートは，線形代数のスペクトル定理とRayleigh商を用いて，図形の最大値を求める．
単なる二次式の平方完成ではなく，許容される向きの制約を持つ幾何学的な最適化を扱う．

\begin{example}[東京大学2025年度・理科数学第3問（要約）]\label{ex:tokyo}
$AB=a,BC=b$，$0<a\le b$，$\angle ABC=\pi/6$ の平行四辺形 $ABCD$ がある．
その各頂点が一つずつ長方形の相異なる四辺上に載るとき，長方形の面積 $S$ を最大化する．
出題では $\theta=\angle BCG$ を用いて面積を表示し，最大値を $a,b$ で記述することを求める．
\end{example}

\begin{example}[京都大学2024年度・理系数学第2問（要約）]\label{ex:kyoto}
$|x|\le2$ と $|y-(8+6i)|=3$ を満たす複素数 $x,y$ に対し，$z=(x+y)/2$ の動く範囲を求め，その面積を計算する．
\end{example}

\section{実対称行列とRayleigh商}
\begin{definition}[二次形式]
実対称行列 $A\in M_n(\mathbb R)$ に対し，$Q_A(x)=x^{\mathsf T}Ax$ を二次形式という．
$x\ne0$ に対する $\mathcal R_A(x)=Q_A(x)/(x^{\mathsf T}x)$ をRayleigh商という．
\end{definition}

\begin{theorem}[Rayleigh--Ritzの原理]\label{thm:rayleigh}
$A$ を実対称 $n$ 次行列とする．その固有値を重複度込みで $\lambda_1\le\cdots\le\lambda_n$ と並べると
\[
 \min_{\norm{x}=1}x^{\mathsf T}Ax=\lambda_1,\qquad
 \max_{\norm{x}=1}x^{\mathsf T}Ax=\lambda_n
\]
である．最大値は最大固有値の固有空間上の単位ベクトルにおいて，最小値は最小固有値の固有空間上の単位ベクトルにおいて，それぞれちょうど達せられる．
\end{theorem}
\begin{proof}
実対称行列のスペクトル定理により，直交行列 $U$ を用いて
$U^{\mathsf T}AU=\mathrm{diag}(\lambda_1,\ldots,\lambda_n)$ と書ける．
$y=U^{\mathsf T}x$ とすると $\sum_jy_j^2=1$ であり，
\[
 x^{\mathsf T}Ax=\sum_{j=1}^n\lambda_jy_j^2
\]
となる．$y_j^2$ は非負で和が $1$ だから，右辺は固有値の凸結合であり，$\lambda_1\le Q_A(x)\le\lambda_n$ である．
固有値の等号条件もこの式から直ちに従う．
\end{proof}

\begin{example}[単位円上の二次式]
$x^2+y^2=1$ のとき $3x^2+4xy$ の最大値は $4$，最小値は $-1$ である．
実際，$A=\begin{pmatrix}3&2\\2&0\end{pmatrix}$ の特性多項式は $\lambda^2-3\lambda-4=(\lambda-4)(\lambda+1)$ であり，定理~\ref{thm:rayleigh} を適用する．
\end{example}

\section{円弧上の制約付き二次形式}
\begin{lemma}[二重角と対称行列]\label{lem:matrix}
$A,B\in\mathbb R$，$R=\sqrt{A^2+B^2}>0$ とする．
\[
 M=\begin{pmatrix}A&B\\B&-A\end{pmatrix},\qquad
 u(\theta)=\begin{pmatrix}\cos\theta\\\sin\theta\end{pmatrix}
\]
とおけば
\[
 u(\theta)^{\mathsf T}Mu(\theta)
 =A\cos2\theta+B\sin2\theta.
\]
$M$ の固有値は $R$ と $-R$ である．
\end{lemma}
\begin{proof}
二次形式を展開して $A(\cos^2\theta-\sin^2\theta)+2B\sin\theta\cos\theta$ とすれば前半が従う．
$\det(M-\lambda I)=\lambda^2-(A^2+B^2)$ なので固有値は $\pm R$ である．
\end{proof}

\begin{proposition}[閉円弧上の最大値]\label{prop:arc}
$A,B>0$，$0<\eta<\pi/2$ とし，$F(t)=A\cos t+B\sin t$ を $0\le t\le\eta$ 上で考える．
$t_0=\arctan(B/A)$ とおくと
\[
 \max_{0\le t\le\eta}F(t)
 =\begin{cases}
 \sqrt{A^2+B^2}&(t_0\le\eta),\\
 A\cos\eta+B\sin\eta&(t_0>\eta).
 \end{cases}
\]
\end{proposition}
\begin{proof}
$R=\sqrt{A^2+B^2}$ とすれば $A=R\cos t_0,B=R\sin t_0$ だから $F(t)=R\cos(t-t_0)$ である．
$0<t_0<\pi/2$ なので $[0,\eta]$ 上の最大値は，$t_0$ が区間内なら $R$，区間外なら右端 $t=\eta$ で得られる．
\end{proof}

\section{平行四辺形と外接長方形}
\begin{lemma}[許容される角と辺の表示]\label{lem:geometry}
例~\ref{ex:tokyo} の条件を満たす長方形について，ある $\theta\in[0,\pi/6]$ が存在し，長方形の横幅 $w$ と高さ $h$ は適当な直交座標系において
\begin{align*}
 w&=a\cos(\theta+\pi/3)+b\sin\theta,\\
 h&=a\sin(\theta+\pi/3)+b\cos\theta
\end{align*}
である．逆にこの区間の任意の $\theta$ に対し，対応する長方形が構成できる．
\end{lemma}
\begin{proof}
長方形の辺を座標軸に平行にとり，$B=(0,0)$ とする．
$BC$ は長さ $b$ で第1象限に向かい，$\angle BCG=\theta$ であるから
$C=(b\sin\theta,b\cos\theta)$ となる．
$BA$ は $BC$ となす角が $\pi/6$ で，左上に向かう．従って
$A=(-a\cos(\theta+\pi/3),a\sin(\theta+\pi/3))$ である．
辺の両端を含むという条件から，横座標および縦座標で必要な射影は非負であり，
$\theta\ge0$ と $\cos(\theta+\pi/3)\ge0$，すなわち $0\le\theta\le\pi/6$ が必要である．
$D=A+C$ なので，四点の各座標の最小値・最大値を比較すると
\[
 x\text{ の範囲}=[-a\cos(\theta+\pi/3),b\sin\theta],
\]
\[
 y\text{ の範囲}=[0,a\sin(\theta+\pi/3)+b\cos\theta]
\]
である．この最小外接長方形は，$A,B,C,D$ が順に左辺・下辺・右辺・上辺に乗る．
元の四辺の命名をこの順序に合わせれば条件を満たす．
逆に任意の $\theta\in[0,\pi/6]$ について上記座標で $A,B,C,D$ を構成すれば，各頂点は対応する辺に乗る．
従って角度条件は必要十分である．
\end{proof}

\begin{proposition}[面積の二次形式表示]\label{prop:area}
補題~\ref{lem:geometry} の設定において
\begin{equation}\label{eq:area}
 2S-ab=\frac{\sqrt3 a^2}{2}\cos2\theta
 +\left(b^2-\frac{a^2}{2}\right)\sin2\theta.
\end{equation}
\end{proposition}
\begin{proof}
$S=wh$ に補題~\ref{lem:geometry} の式を代入して展開する．
交差項は
$ab\{\cos(\theta+\pi/3)\cos\theta+\sin(\theta+\pi/3)\sin\theta\}
=ab\cos(\pi/3)=ab/2$ である．
二乗項はそれぞれ $a^2\sin(2\theta+2\pi/3)/2$ と $b^2\sin2\theta/2$ であり，
$\sin(2\theta+2\pi/3)=\sin(\pi/3-2\theta)$ を用いて整理すると式~\eqref{eq:area} を得る．
\end{proof}

\begin{theorem}[東京大学の図形問題の解答]\label{thm:tokyo}
$0<a\le b$ とする．例~\ref{ex:tokyo} の長方形の面積の最大値は
\begin{equation}\label{eq:answer}
 \max S=\begin{cases}
 \displaystyle\frac{ab+\sqrt{a^4-a^2b^2+b^4}}{2}&(a\le b\le\sqrt2a),\\[4pt]
 \displaystyle\frac{ab}{2}+\frac{\sqrt3 b^2}{4}&(b\ge\sqrt2a)
 \end{cases}
\end{equation}
である．境界 $b=\sqrt2a$ では二つの表示は一致する．
\end{theorem}
\begin{proof}
式~\eqref{eq:area} に対し
$A=\sqrt3a^2/2>0$，$B=b^2-a^2/2>0$，$t=2\theta\in[0,\pi/3]$ とおく．
補題~\ref{lem:matrix} により $2S-ab$ は単位ベクトル上の対称行列の二次形式で，固有値は $\pm R$，
\[
 R=\sqrt{A^2+B^2}=\sqrt{a^4-a^2b^2+b^4}
\]
である．無制約の最大値 $R$ が許される角に現れる条件は，命題~\ref{prop:arc} により
\[
 \arctan(B/A)\le\pi/3
 \iff B\le\sqrt3 A\iff b^2\le2a^2
\]
である．この場合 $S=(ab+R)/2$ となる．
$b^2>2a^2$ では右端 $t=\pi/3$ で最大となり，
\[
 2S-ab=\frac A2+\frac{\sqrt3B}{2}=\frac{\sqrt3b^2}{2}
\]
だから第二式を得る．
$b^2=2a^2$ のとき $R=\sqrt3a^2=\sqrt3b^2/2$ なので二式は一致する．
\end{proof}

\section{Minkowski和によるもう一つの過去問}
\begin{definition}[Minkowski和]
平面の集合 $E,F\subset\mathbb R^2$ に対し，$E+F=\{x+y:x\in E,y\in F\}$ をMinkowski和という．
中心 $0$，半径 $r$ の閉円板を $\overline B_r$，半径 $R$ の円周を $S_R$ と書く．
\end{definition}

\begin{lemma}[円と円板のMinkowski和]\label{lem:annulus}
$R\ge r\ge0$ ならば
\[
 S_R+\overline B_r=\{z\in\mathbb C:R-r\le|z|\le R+r\}.
\]
\end{lemma}
\begin{proof}
$z=u+v$，$|u|=R$，$|v|\le r$ なら三角不等式と逆三角不等式で $R-r\le|z|\le R+r$ となる．
逆に $t=|z|\in[R-r,R+r]$ とする．$z\ne0$ なら $u=Rz/t$ を選ぶと $|u|=R$，$|z-u|=|t-R|\le r$ だから $z=u+v$ と表せる．
$z=0$ がこの範囲に入るのは $R=r$ の場合に限り，この場合は $u$ を任意の長さ $R$ のベクトルとし $v=-u$ とすればよい．
\end{proof}

\begin{theorem}[京都大学の複素数領域問題]\label{thm:kyoto}
例~\ref{ex:kyoto} の複素数 $z$ の動く範囲は
\[
 \frac12\le|z-(4+3i)|\le\frac52
\]
であり，その面積は $6\pi$ である．
\end{theorem}
\begin{proof}
$z-(4+3i)=x/2+(y-(8+6i))/2$ と分解する．
$x/2$ は半径 $1$ の閉円板を，$(y-(8+6i))/2$ は半径 $3/2$ の円周を動く．
補題~\ref{lem:annulus} を $R=3/2,r=1$ に適用すれば領域の表示が従う．
面積は外円の面積から内円の面積を差し引き
$\pi\{(5/2)^2-(1/2)^2\}=6\pi$ となる．
\end{proof}

\begin{remark}[大学数学の役割]
東京大学の問題は許容される円弧上の二次形式の最大化であり，京都大学の問題は集合のMinkowski和の計算である．
いずれも図形の個別の座標計算を，線形代数と集合演算に置き換えている．
\end{remark}

\section*{参考文献・出典}
\begin{enumerate}[label={[\arabic*]}]
\item 東京大学，2025年度前期日程・理科数学，第3問．
\item 京都大学，2024年度前期日程・理系数学，第2問．
\item S.~Lang, \textit{Algebra}, 3rd ed., Springer, 2002, Chap.~XIII (matrices and linear maps).
\item J.~Nocedal and S.~J.~Wright, \textit{Numerical Optimization}, Springer, 1999, Chaps.~2 and 12.
\item 「なかけんの数学ノート」，東京大学2025年第3問・京都大学2024年第2問（出題条件を照合）．
\end{enumerate}

\end{document}